% ======================================================================
\chapter{GitHub Pages e GitHub Actions}
\label{cap:pages}
% ======================================================================

\section{GitHub Pages: um site a partir do repositório}

O \textbf{GitHub Pages} hospeda \textbf{sites estáticos} (HTML, CSS e
JavaScript) direto de um repositório, de graça, no endereço
\texttt{https://<seu-usuario>.github.io/<repositorio>/}
\cite{github_pages}.

\begin{itemize}
  \item Em repositório \textbf{público}, funciona em qualquer conta. Em
        repositório \textbf{privado}, exige o GitHub Pro (que vem com os
        benefícios de estudante, capítulo~\ref{cap:github}).
  \item Aceita domínio próprio, além do \texttt{github.io}.
  \item Não roda código no servidor: nada de PHP, banco de dados ou Python do
        lado do servidor.
\end{itemize}

\subsection{As duas formas de publicar}

A escolha fica em \menu{Settings} \(\rightarrow\) \menu{Pages}
\(\rightarrow\) \menu{Build and deployment} \(\rightarrow\) \menu{Source}:

\begin{table}[H]
  \centering
  \caption{Fontes de publicação do GitHub Pages}
  \label{tab:pages_fontes}
  \begin{tabularx}{\textwidth}{@{}l X@{}}
    \toprule
    \textbf{Source} & \textbf{Quando usar} \\
    \midrule
    \menu{Deploy from a branch} & o site já está pronto no repositório: basta
      escolher a branch (por exemplo, \texttt{main}) e a pasta
      (\texttt{/(root)} ou \texttt{/docs}) \\
    \menu{GitHub Actions} & o site precisa ser gerado antes de publicar, ou
      você quer controlar o que é publicado. É o caminho recomendado pela
      documentação atual \\
    \bottomrule
  \end{tabularx}
\end{table}

\begin{aviso}[O index.html precisa estar na pasta publicada]
O Pages abre o \arq{index.html} da pasta escolhida. Se ele estiver dentro de
uma subpasta (por exemplo, \arq{jogo-da-memoria/index.html}), o endereço
principal responde com erro 404.
\end{aviso}

% ----------------------------------------------------------------------
\section{GitHub Actions: automação dentro do GitHub}
% ----------------------------------------------------------------------

O \textbf{GitHub Actions} executa tarefas automaticamente quando algo acontece
no repositório: rodar testes a cada push, compilar um projeto, publicar um
site \cite{github_actions}.

\begin{table}[H]
  \centering
  \caption{Vocabulário do GitHub Actions}
  \label{tab:actions_termos}
  \begin{tabularx}{\textwidth}{@{}l X@{}}
    \toprule
    \textbf{Termo} & \textbf{O que é} \\
    \midrule
    Workflow & o processo automatizado, descrito num arquivo YAML em
      \arq{.github/workflows/} \\
    Evento (\texttt{on}) & o que dispara o workflow: um \texttt{push}, um Pull
      Request, um horário, um clique manual (\texttt{workflow\_dispatch}) \\
    Job & um conjunto de passos que roda numa mesma máquina \\
    Step & um passo: um comando (\texttt{run}) ou uma ação pronta
      (\texttt{uses}) \\
    Runner & a máquina virtual do GitHub que executa o job, como
      \texttt{ubuntu-latest} \\
    Action & um passo pronto e reutilizável, como
      \texttt{actions/checkout} \\
    \bottomrule
  \end{tabularx}
\end{table}

\begin{dica}[Minutos]
Em repositórios \textbf{públicos}, o GitHub Actions é gratuito. Em privados,
os minutos por mês dependem do plano (tabela~\ref{tab:free_pro}).
\end{dica}

\subsection{Exemplo 1: publicar um site pelo Actions}

Salve como \arq{.github/workflows/pages.yml} e, em \menu{Settings}
\(\rightarrow\) \menu{Pages}, escolha a Source \menu{GitHub Actions}:

\begin{yamlcode}
name: Publicar o site

on:
  push:
    branches: [ main ]
  workflow_dispatch:

permissions:
  contents: read
  pages: write
  id-token: write

concurrency:
  group: pages
  cancel-in-progress: false

jobs:
  publicar:
    runs-on: ubuntu-latest
    environment:
      name: github-pages
      url: ${{ steps.deploy.outputs.page_url }}
    steps:
      - uses: actions/checkout@v7
      - uses: actions/configure-pages@v6
      - uses: actions/upload-pages-artifact@v5
        with:
          path: .
      - id: deploy
        uses: actions/deploy-pages@v5
\end{yamlcode}

O \texttt{path: .} publica a raiz do repositório. Se o site estiver numa
subpasta, troque pelo caminho dela.

\subsection{Exemplo 2: rodar um script e guardar o resultado}

Este workflow roda o \arq{main.py} (disponível em \arq{materials/}), que cria
um arquivo \arq{resultado.txt}, e grava esse arquivo de volta no repositório:

\begin{yamlcode}
name: Executar o script Python

on:
  push:
    branches: [ main ]
  workflow_dispatch:

permissions:
  contents: write   # necessário para o git push no último passo

jobs:
  executar:
    runs-on: ubuntu-latest
    steps:
      # 1. Baixa o repositório para a máquina do GitHub
      - uses: actions/checkout@v7

      # 2. Instala o Python
      - uses: actions/setup-python@v7
        with:
          python-version: '3.x'

      # 3. Roda o script, que cria o resultado.txt
      - name: Executar o main.py
        run: python main.py

      # 4. Grava o resultado no repositório
      - name: Commitar o resultado
        run: |
          git config user.name "github-actions[bot]"
          git config user.email "github-actions[bot]@users.noreply.github.com"
          git add resultado.txt
          git commit -m "chore: atualiza o resultado.txt" || echo "Nada para commitar"
          git push
\end{yamlcode}

\begin{aviso}[Sem permissão, o push falha]
Por padrão, o token que o workflow recebe (\texttt{GITHUB\_TOKEN}) pode ser
só de leitura. Sem o \texttt{permissions: contents: write}, o
\texttt{git push} do último passo falha com erro 403.
\end{aviso}

\subsection{E o \texttt{[skip ci]}?}

Poderia parecer que o push feito pelo próprio workflow dispara o workflow de
novo, num laço infinito. \textbf{Não dispara}: pushes feitos com o
\texttt{GITHUB\_TOKEN} não iniciam novos workflows. O \texttt{[skip ci]} na
mensagem do commit (ou \texttt{[ci skip]}, \texttt{[no ci]},
\texttt{[skip actions]}) só é necessário quando o push é feito com um token
pessoal ou de aplicativo \cite{github_skip}.

\subsection{Acompanhando}

A aba \menu{Actions} do repositório mostra cada execução, com o log de cada
passo. Pelo terminal:

\begin{bashcode}
gh run list          # últimas execuções
gh run watch         # acompanha a execução em andamento
\end{bashcode}
